\documentclass[conference]{IEEEtran}
\usepackage{amsmath}
\usepackage{booktabs}
\usepackage{url}
\begin{document}
\title{Chest X-Ray Pneumonia Classification: Transfer Learning and Feature Fusion}
\author{\IEEEauthorblockN{Project team}\\\IEEEauthorblockA{Replace with team names and institution}}
\maketitle
\begin{abstract}
This educational study reproduces the evaluation of a reference ResNet18 checkpoint and investigates fusion of its learned embeddings with handcrafted image descriptors for binary pneumonia classification. On 624 test images, the reference checkpoint achieved 90.06\% accuracy, 96.67\% sensitivity and 79.06\% specificity. In a logistic-regression feature ablation, fusion achieved 98.46\% sensitivity and 61.11\% specificity, compared with 97.44\% and 69.23\% for CNN embeddings alone. The experiment is exploratory: test images had been inspected previously, patient-level independence is unverified, and no external validation is available. It is not suitable for clinical use.
\end{abstract}
\begin{IEEEkeywords}pneumonia, chest radiography, transfer learning, feature fusion, ablation\end{IEEEkeywords}
\section{Introduction}
Deep learning can learn image representations for binary chest X-ray classification. A meaningful hybrid must be compared with its components under the same classifier and split. We investigate whether coarse handcrafted appearance features complement a learned CNN representation. The trained ResNet18 checkpoint, preprocessing code and manifest were obtained from the reference repository \cite{repo}; the checkpoint was not trained by this project team.
\section{Related Work}
Kermany et al. investigated transfer learning for medical imaging, including pediatric chest radiography \cite{kermany}. He et al. introduced residual connections that ease optimization of deep networks \cite{resnet}. CheXNet demonstrated a deep convolutional approach to chest radiograph classification on a different dataset \cite{chexnet}; its results are not directly comparable to ours. Grad-CAM provides gradient-based visual attribution, not lesion segmentation or proof of clinical reasoning \cite{gradcam}.
\section{Data and Methods}
The source manifest uses the public pediatric chest X-ray dataset \cite{kaggle}. The original training directory is divided into 4,432 training and 784 validation images; the original test directory contributes 624 images (234 normal and 390 pneumonia). Original validation images are excluded. SHA-256 checks against the manifest found no local file mismatches. Patient identifiers are not independently verified. The reference split differs from this repository's earlier group-audited split.

The reference ResNet18 uses grayscale images resized to $224\times224$, repeated into three channels, and normalized using ImageNet statistics. Its ImageNet-initialized backbone was frozen and a dropout/linear binary head was trained with weighted binary cross entropy and AdamW. The reference used random horizontal flips and small rotations in training, validation-loss checkpoint selection and early stopping. The provided checkpoint was selected at epoch nine. During the source training the backbone parameters were frozen, but batch-normalization running statistics can still update in training mode.

For the hybrid, frozen ResNet18 features after global average pooling yield a 512-value embedding. A 68-value handcrafted vector contains an $8\times8$ grid of mean intensities, global mean and standard deviation, and mean absolute horizontal and vertical intensity differences. We compare class-balanced logistic regression on the handcrafted vector, the CNN embedding and their concatenation. Each pipeline fits its scaler and classifier only on the training split. Regularization $C\in\{0.01,0.1,1\}$ is selected using validation balanced accuracy. The same validation data also selected the neural checkpoint, so this is exploratory model development. All test results use a fixed threshold of 0.5.
\section{Evaluation and Results}
Sensitivity is $TP/(TP+FN)$, specificity is $TN/(TN+FP)$, and balanced accuracy is their mean. ROC-AUC summarizes ranking over thresholds. The locally reproduced reference checkpoint produced TN=185, FP=49, FN=13 and TP=377, matching the source confusion counts. Its accuracy was 90.06\%, balanced accuracy 87.86\% and ROC-AUC 0.9551.
\begin{table}[htbp]
\caption{Feature ablation on 624 test images (percent except AUC)}
\centering\small
\begin{tabular}{lrrrr}\toprule
Input & Bal. acc. & Sens. & Spec. & AUC\\\midrule
Handcrafted & 68.63 & 97.95 & 39.32 & 0.9003\\
CNN embedding & 83.33 & 97.44 & 69.23 & 0.9368\\
Fusion & 79.79 & 98.46 & 61.11 & 0.9473\\\bottomrule
\end{tabular}
\end{table}
The hybrid missed six pneumonia cases compared with ten for the CNN-only logistic model, but generated 91 false positives compared with 72. Thus fusion improved sensitivity and AUC in this run while reducing specificity and balanced accuracy. It does not demonstrate an unconditional benefit. Validation balanced accuracy was 94.79\%, 95.31\%, and 96.39\%, respectively, considerably higher than test performance.
\section{Limitations and Conclusion}
The test set has already been inspected and therefore is not a fresh independent confirmation. A single pediatric dataset, unknown patient independence, possible near-duplicates, distribution shift, single checkpoint and no uncertainty intervals limit generalization. The reference scores and earlier custom-CNN report cannot be fairly compared because they use different splits. Future work should verify patient-level splits, repeat experiments with multiple seeds, select operating thresholds on development data and evaluate an untouched external cohort. This prototype is not intended for diagnosis or clinical decisions.
\begin{thebibliography}{00}
\bibitem{kermany} D. S. Kermany et al., ``Identifying medical diagnoses and treatable diseases by image-based deep learning,'' \emph{Cell}, 2018, doi: 10.1016/j.cell.2018.02.010.
\bibitem{resnet} K. He et al., ``Deep residual learning for image recognition,'' CVPR, 2016, doi: 10.1109/CVPR.2016.90.
\bibitem{chexnet} P. Rajpurkar et al., ``CheXNet: Radiologist-level pneumonia detection on chest X-rays with deep learning,'' arXiv:1711.05225, 2017.
\bibitem{gradcam} R. R. Selvaraju et al., ``Grad-CAM: Visual explanations from deep networks via gradient-based localization,'' ICCV, 2017, doi: 10.1109/ICCV.2017.74.
\bibitem{kaggle} P. Mooney, ``Chest X-Ray Images (Pneumonia),'' Kaggle, \url{https://www.kaggle.com/datasets/paultimothymooney/chest-xray-pneumonia}.
\bibitem{repo} Silence91169, ``Pneumonia screening ML/DL,'' commit 814aabe, \url{https://github.com/Silence91169/pneumonia-screening-ml-dl}.
\end{thebibliography}
\end{document}
